\chapter{Discussion}
\label{ch:discussion}

\section{What the results establish}
\label{sec:disc-establish}

Chapter~\ref{ch:results} reports the numbers. This chapter sets out which claims they support, where the pipeline is weakest, how much weight the face/scene contrast can bear, and where the findings stand next to the studies in Chapter~\ref{ch:lit}. The results support three statements.

The first is that the prompt reader's reading of what a visitor typed and the image reader's reading of the drawing made from it agree, on all four rungs, well above the random-pairing null. Clustering on prompt groups discounts the \nRepeats{} records whose prompt appeared earlier, and the intervals in Table~\ref{tab:headline} still clear their nulls. Keeping only the first record of each prompt lowers the rates a little and the nulls with them, and the sign $z$ goes from \FullEmoSignZ{} to \DistEmoSignZ{}. Holm's correction over the \HolmM{} non-degenerate comparisons keeps \HolmSurvive{}. Valence sign carries the most weight here, \FullEmoSignObs{} where the null sits at \FullEmoSignNull{} ($z$ \FullEmoSignZ{}, interval \FullEmoSignCi{}), because it asks the simplest question and states its denominator, \FullEmoNSign{} of \FullEmoN{}.

The second is that the agreement sits in drawings in which the describer found a face. The design set out to test faces against scenes, and the outcome is that only the face drawings agree at all for this reader. Face-classified records reach \FullFaceSignObs{} on valence sign against \FullSceneSignObs{} for scene-classified ones, a gap of \DiffSignPp{} with a clustered interval of \DiffSignCi{} and a permutation $p$ of \DiffSignP{}. On distinct prompts the scene family cell drops to \DistSceneFamObs{} ($z$ \DistSceneFamZ{}) and is among the \HolmFail{} comparisons Holm removes. The template ablation removes the scene-side agreement that is left. With the people-aware template the face rows are identical by construction while the scene rows fall onto their nulls, \AblPeopleAwareSceneSignObs{} against \AblPeopleAwareSceneSignNull{} on sign ($z$ \AblPeopleAwareSceneSignZ{}) and \AblPeopleAwareSceneFamObs{} against \AblPeopleAwareSceneFamNull{} on family ($z$ \AblPeopleAwareSceneFamZ{}).

The third is that the specificity check and the direct-CLIP comparison behave as expected. Where the prompt reader found no affect, the only measure with a working null sits at chance, \CtrlVaObs{} against \CtrlVaNull{} (clustered $z$ \CtrlVaZClustered{}). The control audit flagged \CtrlRemoved{} records, \CtrlFlagged{} containing a plain affect word and \CtrlNameOnly{} reduced by redaction to the bare placeholder. With those removed the clustered $z$ is \CtrlCleanVaZClustered{}. Scored like for like on the \McN{} records, describe-then-extract reaches \McDeFam{} family agreement and direct CLIP \McClipFam{}, which leaves a gap of \McDiffFamPp{} (CI \McDiffFamCi{}, cluster sign-flip $p$ \McClusterFamP{}). There is no ground truth, so the gap does not say which reader is the better one. The two image readers share a top-1 family on only \McReaderAgree{} of these records, and describe-then-extract agrees with the prompt reader far more often.

None of this says what a person would see, because no person has rated the drawings. Every number in Chapter~\ref{ch:results} is one machine reader agreeing with another. The single-coder check is the only step in the project that sets a human reading beside the machine ones, and on its \HcN{} drawings the image reader agrees with the coder's valence sign \HcReaderSign{} of the time (Section~\ref{sec:results-human}). One coder on a sample of that size settles nothing on its own.

\section{The describer is the weak link}
\label{sec:disc-describer}

All of the audits point at the describer. The description template writes a person with a facial expression into every description, whether or not the drawing has a face in it. Of the \DescSceneN{} scene-classified records, the describer counted nobody in \DescScenePeopleZero{}, and for \DescScenePeopleZeroExpr{} of those it still named an expression. For \DescSceneSmile{} of them (\DescSceneSmilePct{} of the no-person records) the word was ``smile'', ``smiling'' or ``happy''. Because the image reader reads only that description, the smile word goes straight into the emotion decision. Where the expression answer is a bare emotion word (\DescBareStrictN{} records, \DescBareStrictPct{} of everything analysed), the reader's top family and the family of that word coincide \DescCopyStrict{} of the time. On those records the image reader mostly returns the family of the describer's own emotion word, so the vision model is taking part in the emotion decision that the two-stage design was meant to keep it out of.

The ablation puts numbers on this (Table~\ref{tab:ablation}). Taking the expression sentence out of the template drops family agreement for the emotional prompts from \FullEmoFamObs{} to \AblDropExpressionEmoFamObs{}, against a null of \AblDropExpressionEmoFamNull{}, and the sign rung from \FullEmoSignObs{} to \AblDropExpressionEmoSignObs{}. With the expression word pinned to ``unclear'' the figure is \AblFixedExpressionEmoFamObs{}. Removing the action, the posture or the scene clause instead moves family agreement to \AblDropActionEmoFamObs{}, \AblDropPostureEmoFamObs{} and \AblDropSceneEmoFamObs{}, a shift of one to three points, so most of the reading rests on the expression clause.

This changes how the comparison with direct CLIP should be read. The gap of \McDiffFamPp{} may simply mean that BLIP-VQA reads a drawn face better than CLIP does, with the two stages adding little, and an ablation that strips emotion words from the expression answer before the reader gets it would tell the two apart. It has not been run. It is listed as further work in Chapter~\ref{ch:conclusion}.

The error taxonomy points the same way (Table~\ref{tab:taxonomy}). Of the \TaxN{} records where both readers were confident and disagreed, the describer accounts for \TaxDescriber{} (\TaxA{} inventing a face or an expression, \TaxB{} misreading one), the image reader for \TaxReader{}, the prompt side for \TaxPrompt{} (\TaxD{} mixed prompts, \TaxF{} prompt-reader misreadings) and the generator for \TaxGenerator{}. So the generator is seldom the cause when the two readings differ, and one of the readers usually is. Bhattacharyya and colleagues \cite{bhattacharyya2025evaluating} find a positive bias in vision-language models. Here it takes the form of a smile rate of \DescSmileFace{} among the \nFaceAll{} face-classified records and \DescSmileAll{} over everything analysed. Gao and colleagues \cite{gao2025eemo} report that the same models mix up the feeling on a depicted face with the feeling a picture is meant to evoke in a viewer, and in this corpus that confusion appears as the \DescScenePeopleZeroExpr{} records where the describer named an expression but counted nobody.

\section{Circularity of the face/scene split}
\label{sec:disc-circularity}

The face/scene contrast was Dr Chen's hypothesis from the meeting at the end of July and the one contrast the design was built to test, so it needs its own caveats. A record is in the face class when the describer counts a person and uses a listed expression word, and in the scene class otherwise. In practice the people count does nearly all of the work, since the describer counted nobody in \DescScenePeopleZero{} of the \DescSceneN{} scene-classified records. That same description is what the image reader then reads. So the contrast compares two kinds of describer output, and its two classes are defined by the same text the reader scores.

The contrast clears its nulls, holds under clustering, and the permutation test gives $p$ = \DiffSignP{} on both sign and family. It is a narrower result than a difference between faces and scenes as drawn, because the classes come from the describer. When a person and an expression are both reported, the family the reader returns tracks the prompt; when nobody is counted, the reader is left with only the expression the template made it assert, and once the people-aware template drops that assertion the scene rows sit on their nulls. What a person would read in a scene drawing is untested. The single-coder check sorts drawings with no input from the describer, and it says the rule is worse than the contrast makes it look: its precision against the coder's judgement of a readable face is \HcFacePrec{}, so most of what it calls a face is not one, and on the coder's own split the face advantage falls to \HcPromptDiffPp{} at $z$ = \HcPromptDiffZ{} (Section~\ref{sec:results-human}). That is a smaller sample and a different pair of readers, so it qualifies the contrast rather than overturning it. None of this was measured until the 9 September audit, though the template dated from early August.

\section{What is not claimed}
\label{sec:disc-not-claimed}

No claim is made that agreement is higher for positive prompts than for negative ones, or the reverse. My first build gave higher agreement for negative prompts than positive ones, and the second build, which had a different image reader, gave the opposite. Both were withdrawn at the audit on 1 September. Inside a single-valence subset every prompt shares one sign, so the valence-sign null has zero spread and there is no test to run, and the raw gap on the valence--arousal rung is a base-rate artefact of where the two families sit on the circumplex. The two image readers also point opposite ways on the direction, so none is reported.

The registered title names what the project set out to study, while the body measures something narrower, the agreement between two machine readers set out in Section~\ref{sec:disc-establish}, and the text says so wherever a number appears because the two are easy to confuse. The human rating study is described as further work in Chapter~\ref{ch:conclusion}.

Colour is outside this study. Every drawing is monochrome and there is no colour condition. Alvarado \cite{alvarado2025texttoimage} find that colour does most of the emotional work in generated images, so a monochrome corpus was expected to be a hard case, but their model uses colour features alone and says nothing about line. So nothing here can be put down to the missing colour as opposed to the readers, the family scheme or the rarity of some families.

The \ConsFam{} within-group agreement against a null of \ConsFamNull{} is a bound on the entire chain, prompt to drawing to reading, so it says nothing about the generator by itself. The image reader's reliability was not estimated, so part of the spread inside a repeat group is measurement noise. Generation parameters are unknown as well. The files record neither seed nor rewritten prompt, so the remainder cannot be split between sampling and prompt rewriting, and a repeat group cannot tell one visitor typing twice from two visitors typing the same word.

Two restrictions remain. Fear (\RecFearN{} prompts), anger (\RecAngerN{}) and disgust (\RecDisgustN{}) occur too seldom among the prompts for any statement at family level, disgust is never returned by the image reader, fear is returned only \famImageFear{} times in \nAnalysed{} records, and surprise is mostly prompts about curiosity. Claims stop at overall agreement and at the face/scene difference. Every number is also a count of records, and the people behind the \nAnalysed{} records number at most \nDistinct{}, so nothing here says how many visitors felt anything.

\section{Relation to prior work}
\label{sec:disc-prior}

Lomas and colleagues \cite{lomas2024improved} come closest. Their secondary result, from 24 raters scoring 240 images for ten emotions, was that an emotion drawn on a person scored higher than the same emotion drawn on a robot. The same direction appears here with two machine readers and prompts that visitors typed themselves, and the ablation then shows that this version of the effect depends on the describer reporting a face, whether or not one is drawn. The two results agree in direction, but because their raters were people and these are machine readers, neither validates the other. WikiArt Emotions \cite{mohammad2018wikiart} found that images with faces get more consistent emotion judgements from annotators, and the face/scene contrast reproduces that, except that the consistency here is a machine's.

The ladder rests on Cowen and Keltner \cite{cowen2017self}. If self-reported experience falls into 27 categories joined by gradients, then a near miss between two readers is information and should be scored as such. Valence--arousal similarity sits at \FullEmoVaObs{} against \FullEmoVaNull{} and top-3 overlap at \FullEmoOvlpObs{} against \FullEmoOvlpNull{} on the same records where family agreement is \FullEmoFamObs{}, so the two readers are near one another far more often than they land on the same family, and a one-to-one score would have discarded that closeness.

The prompt reader bounds everything downstream. The text literature says it will struggle with this register \cite{demszky2020goemotions, yang2023context}. Cohen's kappa for the emotional prompts is \KappaEmo{} (CI \KappaEmoCi{}), in the \KappaEmoBand{} band of Landis and Koch \cite{landis1977measurement}, which is the sort of value human annotators manage on fine-grained text emotion. GoEmotions' own annotators agreed only moderately on many of its 28 categories, and two machine readers with a generator and a drawing in between had no chance of beating that. Carrasco and colleagues \cite{carrasco2024level} find that agreement between people and a generator's labels climbs steeply when the labels are reduced to positive and negative, which is the same pattern as the gap between the sign and family rungs here.

Quoting every value with its null is taken from Zhang and Wan \cite{zhang2025rethinking} and from Lee and colleagues \cite{lee2023holistic}, who find that the usual image--text scores shift when an image is perturbed slightly and track human judgement weakly. A score like that should not be quoted bare. Neither should a similarity whose lowest possible value is 0.412 and which random pairings already push to about \FullNeuVaNull{}. CLIPScore \cite{hessel2021clipscore} is the field's default for this kind of judgement, which made direct CLIP the obvious comparison. It does poorly on these drawings because an embedding trained on captioned photographs has little to say about a monochrome drawn face, and that is the same weakness that led the VQAScore authors \cite{lin2024evaluating} to a question-answering judge. On that view the describer is a VQAScore-style judge aimed at facial expression, which fits how much of the reading rests on its expression answer.

\section{Limitations}
\label{sec:disc-limitations}

\begin{enumerate}
\item Every result is agreement between two machine readers. The only human anchor is the single-coder check, and it is one coder, coding once, over \HcN{} drawings.
\item Which prompts count as emotional is the prompt reader's call, never checked by a person; \CtrlFlaggedPct{} of the control contains a plain affect word, so the control means prompts in which the reader found no affect, which is not the same as prompts without feeling.
\item Circularity affects the face/scene split, because its inputs are the describer's answers, the very text the image reader then reads, and the people count drives it.
\item For every drawing the template writes in a person and an expression. The people-aware template exists only in the ablation, so every main result uses the original one.
\item Low-arousal positive states have no family of their own among the seven, so prompts like ``grateful'', ``calm'' and ``at peace'' all end up as neutral.
\item The family centroids were placed by hand. Means computed per family from the NRC norms \cite{mohammad2018obtaining} would put data in place of that judgement.
\item Each $z$ comes from \nullShuffles{} re-pairings and is an ordering statistic; over seeds 1--5 the headline $z$ for valence sign moves between \ZSeedMin{} and \ZSeedMax{}.
\item Nothing can be regenerated because the generation parameters were never recorded, so the consistency figure cannot be broken down.
\item Redaction removes only the name, so the drawing and the rest of the text stay and some re-identification risk stays with them.
\item Everything was run once, on one machine, so a second describer and a test--retest estimate for each reader are both missing.
\end{enumerate}

\section{Threats to validity}
\label{sec:disc-threats}

The largest threat is to construct validity, because no person has rated the drawings and agreement between two machine readers is all that is measured (Section~\ref{sec:disc-establish}). Two internal threats, the leak from the describer's expression answer into the image reader and the split built from the describer's own answers, were not measured until the 9 September audit (Sections~\ref{sec:results-describer} and \ref{sec:results-ablation}), and the third, pseudo-replication from the \nRepeats{} repeated records, is handled by the clustered inference. External validity is narrow. The corpus is one generator, one festival, mostly one language and one monochrome style, and nothing transfers to colour or to another generator without being measured again. On the statistical side, Holm's procedure covers the \HolmM{} comparisons and the degenerate nulls are now caught from the spread of the null itself. The threat the design could not answer is the human anchor, which is the first item of further work in Chapter~\ref{ch:conclusion}.
